% Generated draft: requires mathematical and source review.
\documentclass[11pt]{article}
\usepackage[margin=25mm]{geometry}
\usepackage{fontspec}
\setmainfont{Latin Modern Roman}
\usepackage{amsmath,amssymb,mathtools,unicode-math}
\setmathfont{Latin Modern Math}
\usepackage{enumitem,needspace,xurl,hyperref}
\setlength{\parindent}{0pt}
\setlength{\parskip}{5pt}
\setlength{\emergencystretch}{3em}
\newcommand{\N}{\mathbb{N}}
\newcommand{\Z}{\mathbb{Z}}
\newcommand{\Q}{\mathbb{Q}}
\newcommand{\R}{\mathbb{R}}
\newcommand{\C}{\mathbb{C}}
\newcommand{\F}{\mathbb{F}}
\newcommand{\A}{\mathbb{A}}
\newcommand{\PP}{\mathbb{P}}
\newcommand{\E}{\mathbb{E}}
\newcommand{\eps}{\varepsilon}
\newcommand{\dd}{\,\mathrm{d}}
\begin{document}
\section*{UnsolvedMath ID: 700990}
\textbf{Source problem code:} catalogue entry
\subsection{Definitions and mathematical statement}
Let \(G\) be an infinite group admitting a non-discrete Hausdorff group topology. A group topology is a topology making multiplication and inversion continuous. Let \(\mathcal L(G)\) be the lattice of all group topologies on the underlying group \(G\), ordered by inclusion, and let \(\mathcal H(G)\subseteq\mathcal L(G)\) consist of the Hausdorff group topologies. Write \(\delta_G\) for the discrete topology. Thus \(\tau<\sigma\) means that \(\sigma\) is strictly finer than \(\tau\).

A finite chain \(\tau_0<\cdots<\tau_r\) has length \(r\). It is non-refinable if, for each \(0\leq i<r\), there is no group topology \(\rho\) satisfying \(\tau_i<\rho<\tau_{i+1}\). When the endpoints are Hausdorff, every intermediate finer topology is Hausdorff, so this definition agrees in \(\mathcal L(G)\) and \(\mathcal H(G)\). An infinite chain means an infinite set of distinct topologies linearly ordered by inclusion. The source does not explicitly define non-refinability for an infinite chain without a specified least element; this qualification affects part (d), as explained below. Positive integer parameters are used in the questions below; the source says ``natural number'' without specifying whether zero is included.

\begin{enumerate}
\item[(a)] For every positive integer \(k\), must there exist a non-refinable chain
\[
\tau_0<\tau_1<\cdots<\tau_k,
\qquad \tau_i\in\mathcal H(G)?
\]
The source reports an affirmative answer for countable nilpotent groups, attributed to A. G. Topale.

\item[(b)] Suppose additionally that \(G\) is nilpotent of class \(k\), where \(k\) is a positive integer. Here nilpotency class \(k\) means that the lower central series, defined by \(\gamma_1(G)=G\) and \(\gamma_{j+1}(G)=[\gamma_j(G),G]\), satisfies \(\gamma_{k+1}(G)=\{1\}\) and \(\gamma_k(G)\neq\{1\}\). For positive integers \(m,n\), suppose that
\[
\tau_0<\tau_1<\cdots<\tau_m,
\qquad
\tau'_0<\tau'_1<\cdots<\tau'_n
\]
are non-refinable chains in \(\mathcal H(G)\) with \(\tau_0=\tau'_0\) and \(\tau_m=\tau'_n\). Must
\[
m\leq nk
\]
hold, and is this bound best possible? The source states that the claim holds for \(k=1\), because \(\mathcal L(G)\) is modular when \(G\) is abelian. Its supplied later discussion labels part (b) solved and reports that the inequality holds but is not sharp, citing V. I. Arnautov (2010). No replacement sharp bound is supplied.

\item[(c)] Does there exist a countably infinite group \(G\) satisfying the opening hypothesis, a positive integer \(k\), a finite non-refinable chain
\[
\tau_0<\tau_1<\cdots<\tau_k
\quad\text{in }\mathcal H(G),
\]
and an infinite chain \(\{\tau'_\gamma:\gamma\in\Gamma\}\subseteq\mathcal L(G)\) such that
\[
\tau_0<\tau'_\gamma<\tau_k
\qquad\text{for every }\gamma\in\Gamma?
\]
The intermediate topologies are automatically Hausdorff because they refine \(\tau_0\). No non-refinability condition is imposed on this infinite chain in part (c).

\item[(d)] Suppose additionally that \(G\) is abelian, and let \(k\) be a positive integer. The intended noncircular reading of the source's definition is
\[
A_k=\bigl\{\tau\in\mathcal H(G):
\text{every non-refinable chain from }\tau\text{ to }\delta_G
\text{ has length }k\bigr\}.
\]
The question is whether, for every infinite non-refinable chain \(C=\{\tau'_\gamma:\gamma\in\Gamma\}\subseteq\mathcal H(G)\) containing \(\delta_G\), one has
\[
A_k\cap C\neq\varnothing.
\]
The source reports that this is true for \(k=1\). Its supplied later discussion labels part (d) solved negatively and reports, more strongly, that no infinite abelian group satisfies this property for \(k=2\), citing D. Peng (2023).

The definition of \(A_k\) and the quantifier over infinite non-refinable chains require clarification for a fully precise formulation: the record does not specify whether chains from \(\tau\) to \(\delta_G\) include infinite chains, whether existence of such chains is required, or the precise non-refinability convention for an infinite chain without a specified lower endpoint. These conventions are left unresolved rather than added as hypotheses.
\end{enumerate}

The opening assumptions have been retained across the subparts. The supplied registry labels the combined problem partially solved. The solution reports and special-case assertions above are supplied source assessments, not independently verified mathematical conclusions.
\subsection{Short English statement}
For an infinite group admitting a non-discrete Hausdorff group topology, the questions concern lengths of non-refinable chains of group topologies, comparisons of lengths for nilpotent groups, coexistence of finite and infinite chains between fixed endpoints, and fixed distances from the discrete topology along infinite chains for abelian groups. The supplied source reports resolutions of (b) and (d). Part (d) needs clarification of its infinite-chain conventions.
\subsection{Sources}
\begin{itemize}
\item Local UnsolvedMath registry, permanent ID \(700990\), exact source code \texttt{catalogue entry}, title \emph{Kourovka Notebook Problem 14.5}; supplied complete source record. Proposed by V. I. Arnautov in 1999. The registry status is \texttt{partially\_solved}.
\item \emph{The Kourovka Notebook: Unsolved Problems in Group Theory}, original problem 14.5 in the 14th issue (1999), as represented in the supplied 21st issue edition (2026), \texttt{arXiv:1401.0300v48}. The registry describes the collection as edited by E. I. Khukhro and V. D. Mazurov.
\item A. G. Topale, deposited in VINITI, 25.12.98, no. 3849--V 98 (Russian). Supplied reference for the affirmative answer to (a) for countable nilpotent groups; no title or link is supplied.
\item V. I. Arnautov, \emph{Bull. Acad. \c{S}tiin\c{t}e Repub. Moldova, Mat.}, 2010, no. 2 (2010), 3--19. Supplied reference for the report that the inequality in (b) holds but is not sharp; no article title or link is supplied.
\item D. Peng, preprint, 2023, \url{https://arxiv.org/abs/2310.08269}. Supplied reference for the negative answer to (d), including the reported failure for every infinite abelian group when \(k=2\); no title is supplied.
\item The supplied status review is dated exactly \texttt{2026-10-16}, with evidence designated \texttt{maintained\_source}, and attributes the solved subparts to the October 2026 21st issue. This is supplied dated curation, not independently verified. Its status-review note is truncated during the Peng reference; the separate supplied background contains the complete link reproduced above.
\end{itemize}
\end{document}
